\section{Introduction}

Heavy work increasingly runs on the web client: image and video processing,
documents, visualization, machine learning. WebAssembly with threads makes it
possible to parallelize that work in the browser~\cite{haas2017}, but it is
not known whether it scales like the same code natively, or whether thread
synchronization widens the gap. A developer who ports a parallel Rust or C++
program to the web needs to know two things: whether adding threads pays as
it does natively, and whether the usual advice about synchronization still
holds. We ask:

\begin{description}[leftmargin=3em]
  \item[RQ1] How do relative speedup and efficiency of the same parallel
  code evolve with the number of threads, natively and in WASM in the
  browser?
  \item[RQ2] To what extent does synchronization between threads affect the
  scalability gap between WASM and native?
\end{description}

\paragraph{Contributions.}
\begin{enumerate}[leftmargin=2em]
  \item A controlled comparison, with the analysis plan fixed before the
  data, of one Rust/Rayon source compiled with the same compiler and release
  profile for native and WASM, in two engines (V8 in Chrome, SpiderMonkey in
  Firefox), with the run as the unit of analysis.
  \item A synchronization design that separates the cost of shared atomics
  from the effect of the input, swept over four contention levels.
  \item An explanation of each finding, tested against predictions written
  before the data: engine arms that select one compiler, a native probe that
  changes one instruction sequence, hardware counters of the loop alone, and
  the machine code each engine runs.
  \item The problems of browser benchmarking that we met and how each was
  detected and controlled.
\end{enumerate}

\paragraph{The answers in one paragraph.} RQ1: with enough work per parallel
region, WASM threads scale like native code and cost a constant factor in
time; the browsers' fork-join is several times slower, which matters only
for small regions. RQ2: the effect of synchronization is a property of the
engine's compiler, not of WASM threads. Chrome pays more than native under
concentrated contention and less under spread contention, because of one
instruction sequence; Firefox pays what native code pays.

\paragraph{How to read this report.} It is self-contained: no claim needs the
repository to be understood. Sections~\ref{sec:background}
and~\ref{sec:related} give the background and the related work.
Sections~\ref{sec:overview} to~\ref{sec:plan} are the method: the study in
one picture and its three builds, the machine, the builds, the workloads,
the harness and the layout of a run, the pilot decisions and the analysis
plan. Section~\ref{sec:quality} is the data quality of every campaign.
Sections~\ref{sec:rq1} and~\ref{sec:rq2} are the primary results,
Section~\ref{sec:why} what explains them, and Section~\ref{sec:exploratory}
the exploratory campaigns. Sections~\ref{sec:problems}
to~\ref{sec:discussion} give the problems found, the threats to validity,
how to reproduce the study, the discussion and what stays open.
Appendix~\ref{sec:appendix} lists every campaign
(Table~\ref{tab:campaigns}) and maps each section to the decision records
and tables behind it (Table~\ref{tab:sources}).

Three levels of claim are kept apart throughout. A \emph{measured result}
is a number from the committed data with its interval. An explanation is
\emph{answered} when an intervention that changes one thing reproduces the
effect, \emph{supported} when independent observations agree with a
prediction written before them, and \emph{open} otherwise. Only the ten
contrasts of the analysis plan are primary; everything else is labelled
exploratory and reported with uncorrected 95\% intervals.

Every result in this report comes from one build, tagged
\code{measured/v3}. Numbers from pilots and from earlier builds appear only
as the evidence for a design decision and are labelled as such.
